\section{Tensor prices and wire distances}\label{sec:prices-background}
A bilinear algorithm of rank \(r\) constructs its tensor from \(r\)
independent scalar products by local linear maps. A \emph{spectral
point} assigns a value to tensors that can only decrease under such
maps. Tensoring multiplies the values, so their logarithms add. We call
these logarithms \emph{prices}. Applied to a wire, the price measures
the cost of sharing an index between its two sites. This is how a
spectral point gives a geometry on the sites of a tensor network.

\subsection{From tensors to distances}
Fix a field and a finite set \(S\) of at least two sites. Each tensor
has one finite-dimensional vector space at each site. A tensor \(T\)
is a \emph{restriction} of \(U\) if one linear map at each site takes
\(U\) to \(T\). For a tensor product, group the factors at each site.
For a direct sum, take the direct sum of the spaces at every site, so
the two tensors occupy separate blocks. The unit \(\mathbf 1\) is the
tensor \(1\) in one-dimensional spaces.

A spectral point \(\varphi\) is a restriction-monotone map to
\(\mathbb R_{\ge0}\) that preserves zero and the unit and satisfies
\[
     \varphi(T\oplus U)=\varphi(T)+\varphi(U),\qquad
     \varphi(T\otimes U)=\varphi(T)\varphi(U).
\]
Every nonzero tensor restricts to \(\mathbf 1\), so its value is at least
one. The price at scale \(D>0\) is \(H(T)=D\log\varphi(T)\), with
\(H(0)=-\infty\). In this section and Part~I logarithms are natural.
Thus the prices of~\cite{oaG} are multiplied by \(\log 2\) at the same scale.

For a nonempty \(A\subseteq S\), the tensor
\[
         \Delta_A(q)=\sum_{j=1}^q\bigotimes_{i\in A} e_j
\]
shares a \(q\)-valued index among the sites of \(A\). Its other legs are
one-dimensional. Tensoring diagonals multiplies their dimensions, and a
local projection can discard any set of labels. Comparing powers of
dimensions \(q\) and \(2\) therefore gives
\[
            H(\Delta_A(q))=\delta(A)\log q,\qquad q\ge2.
\]
We call \(\delta(A)\) the common-label rate on \(A\). In particular
\(\Delta_S(q)\) is a direct sum of \(q\) units, so \(\delta(S)=D\).
For \(a\ne b\), \(\Delta_{\{a,b\}}(q)\) is a wire of dimension \(q\).
Write \(d(a,b)=\delta(\{a,b\})\) and set \(d(a,a)=0\).
For example, matrix rank across a fixed cut of \(S\) is a spectral
point. At scale one its wire distance is one for sites on opposite
sides of the cut and zero for sites on the same side.

Joining two wires by contracting their common endpoint gives a wire
between their other endpoints. Monotonicity and additivity of price
give the triangle inequality for \(d\). A \(q_i\)-dimensional wire
can also move a leg of dimension \(q_i\) from \(x_i\) to \(y_i\).
Moving each leg of a nonzero tensor \(T\), then applying the same
argument in reverse, gives the transport bound
\begin{equation}\label{eq:prices-transport}
       |H(T;x)-H(T;y)|\le\sum_i d(x_i,y_i)\log q_i .
\end{equation}
Here \(x=(x_i)\) and \(y=(y_i)\) are the placements of the legs. Group
the legs assigned to one site. When all legs are at a single site, any
nonzero tensor there is equivalent to the unit and has price zero.

\subsection{Geometric extension and spectral comparison}
A \emph{spectral space} is a metric space with compatible prices for
all finite placements, whose wire rates are the given distances.
Compatibility allows trivial legs to be added and legs at coincident
sites to be grouped. On a finite set of positive diameter the scale is
its common-label rate, which may vary with the set. At a single point the
scale is zero. Identifying sites at zero wire distance turns a spectral
price on finitely many sites into a spectral space.

Every spectral space embeds, with all its prices, in a complete
\(\operatorname{CAT}(0)\) spectral space. In the extension every pair
\(a,b\) has a unique midpoint \(m\), and every point \(o\) satisfies
\[
      d(o,m)^2\le \tfrac12 d(o,a)^2+\tfrac12 d(o,b)^2
                                      -\tfrac14 d(a,b)^2 .
\]
Prices are jointly convex as their legs move along geodesics. Distances
and prices may be rescaled together. In a pointed metric ultralimit
of these spaces, the price of a fixed tensor on a bounded sequence of
placements is the ultralimit of its prices
\citep[Sections~\ext{G}{sec:spectra} and~\ext{G}{sec:limits},
Theorem~\ext{G}{thm:geometry} and Proposition~\ext{G}{prop:joint-convexity}]{oaG}.
We also use convexity of distance and nonexpansiveness of nearest-point
projection onto a nonempty closed convex set. A finitely supported
probability measure has a unique barycenter \(b\) minimizing its mean
squared distance. Evaluating the same mean at \(z\) increases it by at
least \(d(z,b)^2\) \citep{BridsonHaefliger,oaG}.

Over an algebraically closed field, price sublevel sets in a fixed tensor
format are Zariski closed. Under a field extension, the possible joint
values of spectral points on tensors defined over the original field
are preserved \citep[Section~\ext{G}{sec:spectra}]{oaG}.
Finally, if two nonzero tensors on the same sites satisfy
\(\varphi(T)<\varphi(U)\) for every spectral point, there is a fixed
integer \(C\) such that, for every \(j\ge1\), \(T^{\otimes j}\)
is a restriction of the direct sum of \(C\) copies of \(U^{\otimes j}\).
We use this strict form of spectral comparison to turn price inequalities
into exact tensor restrictions \citep[Lemma~\ext{G}{lem:spectral-comparison}]{oaG}.
